\documentclass[11pt,a4paper]{article}

\usepackage[utf8]{inputenc}
\usepackage[T1]{fontenc}
\usepackage[scale=0.75]{geometry}
\usepackage{microtype}
\usepackage{hyperref}

\renewcommand{\rmdefault}{ppl}
\renewcommand{\sfdefault}{phv}
\microtypesetup{expansion=false}
\setlength{\parindent}{0pt}
\setlength{\parskip}{0.6em}

\begin{document}

\section*{Application Text}

I am applying for the PhD position because it combines rigorous machine learning research with real system constraints. My MSc in Statistical Learning and Data Analytics at KTH, together with my Engineering Physics background, gives me the quantitative preparation for doctoral work on trustworthy AI.

At RaySearch, I worked on automated radiotherapy planning using deep learning, probabilistic modelling, and multi-objective optimisation. My degree project was on \href{https://kth.diva-portal.org/smash/get/diva2:1431327/FULLTEXT02.pdf}{Image Distance Learning for Probabilistic Dose-Volume Histogram and Spatial Dose Prediction in Radiation Therapy Treatment Planning}. That work reinforced careful problem formulation and close collaboration with domain specialists.

More recently, I have delivered production ML systems at Valcon and Signum, including real-time pipelines and embedding-based search workflows. I have also worked with non-technical stakeholders to keep outputs actionable. I would bring that mix of research discipline, implementation experience, and clear communication to the lab.

\end{document}
